\label{ch:conclusiones}
Las conclusiones resumirán la viabilidad, precisión e impacto administrativo del agente inteligente propuesto para optimizar los procesos de conciliación contable de la UPCH. Se discutirá cómo la incorporación de lógica basada en modelos de lenguaje grandes resuelve el desafío de la discrepancia semántica y el lenguaje informal del registro clínico, superando las limitaciones tradicionales de la búsqueda por palabras clave o similitud simple.

Se abordará también el valor operativo de la interfaz con supervisión humana (\textit{Human-in-the-Loop}), que permite aliviar la sobrecarga administrativa diaria de la clínica de manera segura sin delegar por completo la responsabilidad legal contable de la facturación. Finalmente, se plantearán las limitaciones identificadas en el estudio y se ofrecerán recomendaciones clave para la integración definitiva del sistema de agentes con las bases de datos transaccionales activas de la institución universitaria.
